\section{Data attribution: Notation and problem setup}
\label{sec:problem_setup}
The high-level goal of data attribution is to connect the choice of training
data to model behavior. For example, one may want to use data
attribution to find the training datapoints that cause a given output, or to
surface data that harms accuracy. In this section, we formalize this
goal with the {\em predictive data attribution} (or
\textit{datamodeling}~\citep{ilyas2022datamodels}) framework, which phrases data
attribution as the task of predicting how model behavior changes as a function
of the training data.

Specifically, we view the machine learning pipeline as a three-step process
wherein we
(a) choose training data;
(b) apply a learning algorithm to that data, yielding a trained model;
and then (c) evaluate the trained model.
The goal of predictive data attribution is to construct a function
that {\em directly} predicts the output of step (c) from the choice of
training data in step (a).
To make this more precise (and borrowing from~\citep{ilyas2022datamodels}),
we define the following notation:
\begin{itemize}
  \item Let $\trainset = \{z_i\}_{i=1}^n$ be a pool of $n$ possible training
    examples. We represent {\em datasets} as vectors $\wvec \in \mathbb{R}^n$
    where each entry $w_i$ is an importance weight for the $i$-th example in
    $\trainset$. The importance weight $w_i$ controls the scaling of the loss of
    sample $i$; for example, $w_i=0$ implies that we do not include the $i$-th
    example in the training set, while $w_i > 0$ implies that we do include the
    example (and multiply its loss by $w_i$ during training).
  \item Let $\alg{}: \mathbb{R}^n \to \Theta$ be a {\em learning algorithm}
    mapping datasets---parameterized by importance weight vectors
    $\wvec$---to trained model
    parameters.
    We assume that all aspects of the training setup beyond the
    training data
    are captured by $\alg{}$ (e.g., learning rate, weight decay, etc.).
  \item Let $\phi: \Theta \to \mathbb{R}$ be a {\em measurement
    function} mapping a machine learning model $\theta$ to a scalar measurement
    $\phi(\theta) \in \mathbb{R}$. For example, $\phi(\theta)$ might represent
    the loss of the classifier with parameters $\theta$ on a given test sample.
  \item Let $f: \mathbb{R}^n \to \mathbb{R}$ be the {\em model output function}
    $f$ mapping datasets directly to model outputs, i.e., a composition
    of $\alg{}$ and $\phi$.
\end{itemize}
To illustrate this notation, we instantiate it in the context of linear
regression. In this case, the training pool is a set of $n$ input-label pairs
$\trainset = \{(\mathbf{x}_i \in \mathbb{R}^d, y_i \in \mathbb{R}) \}_{i=1}^n$;
the learning algorithm $\alg{}$ fits a linear model minimizing the average
squared loss {\em weighted} by a given $\wvec \in \mathbb{R}^n$, and the
measurement function $\phi$ evaluates a model's loss on a specific test point
$\mathbf{x}_{test}$, i.e.,
$$\alg{}(\wvec) := \arg\min_\theta \sum_i w_i \cdot (\theta^\top \mathbf{x}_i -
y_i)^2 \qquad \text{and} \qquad \phi(\theta) := (\theta^\top \mathbf{x}_{test} -
y_{test})^2.$$ Then, $f(\wvec) := \phi \circ \alg{}$ maps a data
weighting $\wvec$
to the resulting model's loss on $\mathbf{x}_{test}$.
Now, in this linear regression example, $f(\wvec)$ is easy to directly compute
(in fact, it has a closed form in terms of $\wvec$),
but this is seldom the case.
In general, evaluating $f(\wvec)$ requires re-training a model on the weighting
vector $\wvec$---which can make be very expensive for large-scale models.
This motivates our (informal) definition of predictive data attribution,
given in Definition~\ref{def:predictive_data_attribution} below.

\begin{boxdef}[Predictive data attribution]
  \label{def:predictive_data_attribution}
  \vspace*{-0.5em}
  A \underline{\smash{predictive data attribution}} is an
  \underline{\smash{explicit}} function
  $\hat{f}$ that aims to approximate a model output function $f$.
  In other words, for a given data weight vector
  $\wvec$, $\hat{f}(\wvec)$ should be fast to compute
  while also accurately predicting $f$, i.e., satisfying
  $\hat{f}(\wvec) \approx f(\wvec)$.
  \vspace*{0.4em}
\end{boxdef}

\subsection{Single-model predictive data attribution}
\label{sec:single_model_data_attribution}
The main goal of this work is to operationalize predictive data attribution
for large-scale (deep) learning algorithms.
A challenge, however, is that these learning algorithms are {\em
non-deterministic},
meaning that the same training dataset can map to \textit{different} model
parameters depending on randomness (e.g., random parameter
initialization or data order shuffling) \citep{zhuang2022randomness,jordan2024variance}.

As a result, when learning algorithms are non-deterministic, we cannot perfectly
predict how choice of data will change model behavior; we can only predict how
it will change the {\em expected} behavior. More precisely, data attribution
methods predict how (on average, over training randomness) a {\em new} model
would behave if retrained from scratch, {\em and not how a specific trained
model would behave if we had changed the training dataset}.

To make this more precise, let $\wvec$ be a weighting vector defining a 
dataset, and let $\hat{f}$ be a data attribution method.
Now, consider the expected difference between our estimator $\hat{f}(\wvec)$,
and the true model output $f(\wvec)$, averaged over training randomness. 
This error decomposes into two~terms:
\begin{equation}
  \label{eq:bias_variance}
  \mathbb{E}\left[(\hat{f}(\wvec) - f(\wvec))^2\right] =
  \underbrace{\left(\hat{f}(\wvec) -
  \mathbb{E}[f(\wvec)]\right)^2}_{\text{Reducible Error}} +
  \underbrace{\mathbb{E}\left[\left(f(\wvec) -
  \mathbb{E}[f(\wvec)]\right)^2\right]}_{\text{Irreducible Error}}
\end{equation}
Looking at each term in \eqref{eq:bias_variance}: the \textit{reducible error}
(or bias) is minimal when $\hat{f}(\wvec) = \mathbb{E}[f(\wvec)]$,
while the irreducible error (or variance) depends only on $f$,
and is constant regardless of the data attribution method $\hat{f}$.
Indeed, the irreducible error
arises from inherent randomness in the model training process, and
thus is fundamentally {\em unattributable} to data. Accordingly, current data
attribution methods can answer questions about algorithms (e.g., {\em ``what
    would happen if we trained a new model on a dataset not containing
    the training
example $x$?''})---but not about individual models.

However, in practice we often want to answer questions about
\textit{individual} models,
not a class of learning algorithms. For example, we might ask a question
like {\em ``what was the effect of training example $x$ on this specific
model?''} 
This question motivates us to define and consider a problem 
that we call {\em single-model data attribution}.




\medskip
\noindent\textbf{Single-model predictive data attribution.}
To understand how choice of data changes \textit{individual} trained models, we
propose a new setting called {\em single-model data attribution}.
Here, we enforce that
the learning algorithm $\alg{}$ and the measurement function $\phi$ are
deterministic (i.e., by fixing data ordering,
parameter initialization, etc.).
This determinism ensures that for any weighting $\wvec$, the model output
$\phi(\alg{}(\wvec))$ is deterministic
(i.e., so that the expected model output that datamodels predict is constant,
and $\text{Var}\left[{\phi(\alg{}(\wvec))}\right] = 0$). In this new setting,
model outputs vary only from changes in training data weights, allowing for
predictive data attribution methods to exactly attribute changes to
data weights. In the language of \eqref{eq:bias_variance}, the irreducible error
is zero.

\begin{remark}[Single-model versus standard predictive data attribution]
  Our motivation for the single-model setting is that in many cases
  we often want to understand the effect of training data on a
  specific model, rather than a class of models.
  Still, even in cases where we \underline{do} care more about 
  the average behavior of a learning algorithm (and not  
  a specific trained model), a near-perfect single-model data attribution
  method can be used to construct a near-perfect standard data
  attribution method by averaging over different learning algorithm seeds.
  We discuss this connection further in Section \ref{sec:discussion}.
\end{remark}



